\section{A compact model and its growth operator}\label{sec:operator}

This section replaces the boundary-growth problem by growth of a
single positive operator on a compact family of normalized maps.

\subsection{Univalence estimates}

We use the classical quarter and derivative-distortion theorems in
exactly the following form \cite[Chapter 2]{Duren}.

\begin{lemma}[Classical univalence estimates]\label{lem:classical}
For every \(f\in\SSS\) and \(|z|=r<1\),
\begin{equation}\label{eq:koebe}
 \frac{1-r}{(1+r)^3}\le |f'(z)|\le
 \frac{1+r}{(1-r)^3},\qquad
 B(0,1/4)\subset f(\DD).
\end{equation}
If \(F\) is univalent on the Euclidean disk \(B(z_0,s)\), then
\begin{equation}\label{eq:local-quarter}
 B\bigl(F(z_0),s|F'(z_0)|/4\bigr)
 \subset F(B(z_0,s)).
\end{equation}
\end{lemma}
\begin{proof}
The first assertions are the quoted classical estimates for normalized
univalent disk maps. Apply the quarter theorem to
\((F(z_0+s\zeta)-F(z_0))/(sF'(z_0))\) to obtain the last assertion.
\end{proof}

Let \(\HH=\{x+iy:y>0\}\), and define
\[
 \KK=\{F:\HH\to\C\text{ univalent}:F(i)=0,\ F'(i)=1\}.
\]
For \(F\in\KK\), set \(Q_F=1/F'\).
With \(m(z)=(z-i)/(z+i)\), the maps of \(\KK\) are exactly
\(F=2i f\circ m\), \(f\in\SSS\).

\begin{lemma}[Compact normalized half-plane family]\label{lem:compact}
The locally uniform topology makes \(\KK\) compact metrizable.
Evaluation of \(F,Q_F\) and all their derivatives is jointly
continuous in \(F,z\), and uniformly bounded on each compact
source set. Moreover, uniformly in \(F\in\KK\),
\begin{equation}\label{eq:strip}
 |Q_F(x+iy)|\le C/y\quad(|x|\le2,\ 0<y\le1),
 \qquad |F(z)|\ge c>0\quad(|m(z)|\ge1/2).
\end{equation}
\end{lemma}
\begin{proof}
Integrating the derivative bound in the disk gives local uniform
boundedness. Montel's theorem gives convergent subsequences.
Limits retain the normalization by convergence of derivatives
and remain injective by Hurwitz's theorem: for a fixed source point,
subtract its value and apply Hurwitz on the punctured domain.
The nonconstant normalization excludes the identically zero
alternative. A countable compact exhaustion metrizes this topology.
Interior holomorphic convergence, together with nonvanishing
derivatives, gives the asserted continuity and bounds.

The identities
\[
 1-|m(z)|^2=\frac{4y}{|z+i|^2},
 \qquad |Q_F(z)|=\frac{|z+i|^2}{4|f'(m(z))|}
\]
and \eqref{eq:koebe} imply the strip estimate.
The image under \(f\) of the disk of radius \(1/2\) contains
\(B(0,1/8)\) by \eqref{eq:local-quarter}. Injectivity gives
the second assertion, with \(c=1/4\) for \(F=2i f\circ m\).
\end{proof}

\subsection{Affine rebasings and dyadic samples}

Identify a point \(z=x+iy\in\HH\) with the affine map
\[
 z\cdot w=x+yw,\qquad z^\#=(-x+i)/y.
\]
This is a group with identity \(i\) and inverse \(z^\#\).
For \(F\in\KK\), define
\[
 (T_zF)(w)=\frac{F(z\cdot w)-F(z)}{yF'(z)}.
\]
Then \(T_zF\in\KK\), and direct substitution gives
\begin{equation}\label{eq:action}
 T_w(T_zF)=T_{z\cdot w}F,\qquad
 Q_{T_zF}(w)=\frac{Q_F(z\cdot w)}{Q_F(z)}.
\end{equation}
On \(C(\KK)\), with its supremum norm, put
\[
 (A_z\Phi)(F)=|Q_F(z)|\,\Phi(T_zF),\qquad
 L=\tfrac12(A_{-1/2+i/2}+A_{1/2+i/2}).
\]
These are bounded positive operators. For each continuous \(\Phi\),
\(A_z\Phi\) is jointly continuous in \(F,z\); rebasings send compact
source sets to compact sets locally in the parameter \(z\).
Equation~\eqref{eq:action} gives \(A_zA_w=A_{z\cdot w}\).
In particular, if
\(v=2^{-n}\) and \(u_j=-1+(2j+1)v\), then
\begin{equation}\label{eq:dyadic}
 L^n=v\sum_{j=0}^{2^n-1}A_{u_j+iv},\qquad
 \|L^n\|=\max_{F\in\KK}(L^n1)(F).
\end{equation}
The abscissae are the sums \(\sum_{\ell=1}^n\pm2^{-\ell}\);
positivity proves the norm identity.
Testing the affine map \(F(z)=z-i\), and using \eqref{eq:strip},
we obtain \(1\le\|L^n\|\le C2^n\).
The same operator will control every map. Submultiplicativity defines
\begin{equation}\label{eq:rho}
 \rho=\lim_{n\to\infty}\|L^n\|^{1/n},
 \qquad \beta=\log_2\rho,\qquad 0\le\beta\le1.
\end{equation}
Indeed, writing \(n=qN+r\), \(0\le r<N\), bounds the upper
root limit by \(\|L^N\|^{1/N}\); take the infimum over \(N\).

\begin{proposition}[From dyadic samples to disk means]\label{prop:disk}
There is a universal constant \(C_0\) such that, for \(f\in\SSS\),
\[
 M_{-1}[f'](r)\le C_0\|L^n\|
 \quad\text{if}\quad 1/2\le r<1,\quad
 v=2^{-n}\le1-r<2v.
\]
Consequently \(B_b(-1)\le\beta\).
\end{proposition}
\begin{proof}
For \(3\pi/4\le\theta\le5\pi/4\), write
\[
 m^{-1}(re^{i\theta})=x+iy,\qquad
 x=\frac{-2r\sin\theta}{d},\quad
 y=\frac{1-r^2}{d},\quad
 d=1-2r\cos\theta+r^2.
\]
Here \(|x|\le1\), \(1/2\le y/v\le4\), and
\[
 \frac{dx}{d\theta}
 =\frac{4r^2-2r(1+r^2)\cos\theta}{d^2}\ge\frac1{16}.
\]
If \(|x-u_j|\le v\), the point
\((x-u_j)/v+i y/v\) lies in a fixed compact rectangle in \(\HH\).
Apply Lemma~\ref{lem:compact} to \(Q_{T_{u_j+iv}F}\) and
\eqref{eq:action} to get
\[
 |Q_F(x+iy)|\le C|Q_F(u_j+iv)|.
\]
For \(F=2i f\circ m\),
\(|f'(re^{i\theta})|^{-1}=4|Q_F(x+iy)|/|x+i(y+1)|^2\).
The denominator is bounded below. Integrating over the cells gives
a bound \(Cv\sum_j|Q_F(u_j+iv)|\le C\|L^n\|\) for this arc.
Four arcs obtained by applying the same argument to
\(e^{-ia}f(e^{ia}\zeta)\) cover the circle.
All constants are independent of \(f\).
Finally \(n/\log(1/(1-r))\to1/\log2\); use \eqref{eq:rho},
take each map's limit superior, and then take the supremum.
\end{proof}

